Include a header paragraph specific to your product here. Functional requirements are the directly observable features and functions of the product encountered by its users --- the ``look and feel'' and behavior that each end user should expect the product to do and/or not do. These correspond to the customer-facing requirements agreed with the intended customer/user/sponsor, and must not be changed without their specific agreement. Each requirement below is written as a singular, verifiable ``shall'' statement and carries the attribute set defined in Section~2.

\subsection{FR-01: [Requirement Name]}
\subsubsection{Description}
A detailed, singular, and verifiable statement of the feature/function. For example: \textit{The GUI background shall be slate blue (\#007FFF, six-digit hexadecimal), so that the GUI matches other software products offered by the customer.} It is acceptable and advisable to include drawings/graphics if they aid understanding.
\subsubsection{Rationale}
The customer need or benefit that this requirement satisfies. This justifies the requirement and supports future change decisions.
\subsubsection{Source}
The origin of the requirement (customer, sponsor, named team member, federal regulation, CSE Senior Design project specifications, etc.).
\subsubsection{Priority}
Critical / High / Moderate / Low / Future
\subsubsection{Verification Method}
Inspection / Analysis / Demonstration / Test --- how conformance will be shown.
\subsubsection{Constraints}
Realistic constraints relevant to this requirement (economic, environmental, social, political, ethical, health \& safety, manufacturability, sustainability), as appropriate.
\subsubsection{Applicable Standards}
Any specific standards that apply (e.g. a color specification, a data-format RFC, an accessibility guideline). Research and document which standards will be followed.

\subsection{FR-02: [Requirement Name]}
Repeat the attribute block for each functional requirement.
